% !TeX root = ../papers/nc-p3.tex
\section{五次対称行列の過剰曲線と二次正規化の全排除}
\label{sec:quintic-web-curve-excess}
\begin{theorem}\label{thm:quintic-web-curve-excess}
被約四辺有限標準Frobenius窓の一配置 $X$ を含む五行対称一次行列の階数三以下スキームを $Z$ とする。
階数二以下の点がなく、$Z$ の正次元部分が高々曲線、その被約総次数が五以下なら、
$Z=X$ である。この結論は一般点で非被約な過剰曲線にも適用できる。
従って被約四辺有限標準正則窓の非恒等反対根の余核は、
二次曲面の二次系による有限双有理射 $\nu:Q\to D$ の $\nu_*\mathcal O_Q(1)$ にはなれない。
この排除に余因子曲線や階数一点の有限性を仮定しない。
\end{theorem}
\begin{proof}[局所自由な外積と次数三での消滅]
$Z$ に曲線があるとし、最大有限長部分を構造層から割った純Cohen--Macaulay曲線を $C$ とする。
スキーム次数を $d$、有限長部分の長さを $a$ とする。五行行列 $A$ の余核を $M$ とし、
$W=(\bigwedge^2 M)/\text{有限長部分}$ と書く。
定理\ref{thm:cofactor-curve-duality}の複体の証明を次数 $(4,5,6)$、階数 $(15,24,10)$ で適用すると
\[
         \chi(\mathcal O_Z(t))+\chi(W(2-t))=20,\qquad H^1(W(-1))=0 .
\]
実際左二項の完全性は余次元二の階数条件による。余因子への写像の第一ホモロジー $T$ は純一次元で、
最左の交代行列項を $\mathcal O(-4)$ に双対化すると
$\mathcal Ext^2(T,\mathcal O(-4))=W(2)$ となる。
$T(3)$ は自由分解 $\mathcal O(-3)^{10}\to\mathcal O(-2)^{24}\to\mathcal O(-1)^{15}$ の余核の部分層なので
$H^0(T(3))=0$、双対で後の消滅を得る。前の等式の定数は元の複体のEuler多項式から二十である。

三行を消去した残りは対称二行で、その全成分が $Z$ を定義する。従って $\bigwedge^2 M$ は
$Z$ 上の可逆層であり、$W$ はその $C$ への制限である。
実際 $A:\mathcal O^5\to\mathcal O(1)^5$ の $Z$ 上の核を $K$、$\mathcal O^5/K$ を $E$ とすると、
$E$ は階数三で $E=E^*(1)$、$M|_Z=K^*$。ゆえに
$W=\det E|_C$, $W^2=\mathcal O_C(3)$ であり、これは各被約成分へも制限できる。
特に $d$ と各成分次数は偶数である。$L=W(-1)$ とすると $L^2=\mathcal O_C(1)$ であり、
上の消滅と一次元のRiemann--Rochにより
$\chi(\mathcal O_C)=h^0(L)-d/2$。非被約でも直線束の次数は基本サイクル上の次数で加法的である。
$X$ のうち曲線の外の点は各々少なくとも局所長一を $a$ に寄与するので
\[
 |X\cap C_{\rm red}|\ \ge\ 20-a
   =\tfrac72 d+2\chi(\mathcal O_C)=\tfrac52 d+2h^0(L).
\]

\end{proof}
\begin{proof}[各台と肥厚の比較]
各成分の次数が偶数なので被約台は一つの円錐曲線、一つの整四次曲線、または二つの円錐曲線である。
円錐曲線上の印は三次評価の全射性から高々七、整四次曲線上は高々十三
（正規化上で次数十二の束の切断は高々十三）、二円錐曲線上は高々十四である。
肥厚があればスキーム次数は被約次数より少なくとも二大きく、$5d/2$ が各上界を超える。
よって $C$ は被約である（純性により一般被約なら全体で被約）。

円錐曲線で $\chi=1$ なら $7d/2+2\chi=9$ で矛盾する。
平面四次では二円錐曲線の和を含め $\chi=-2$ だが、必要な十印は前補題の平面上界九に反する。
非平面整四次は種数上界一で $\chi\ge0$、必要な印数十四以上は十三の上界に反する。
非共面の二円錐曲線の交長を $r$ とすれば $r\le2$、$\chi=2-r$。
各円錐曲線上の三次式の制限はこの交叉上で全射なので、
和の上で $h^0(\mathcal O_C(3))=14-r=12+\chi$ である。
これは必要な $14+2\chi$ に足りない。曲線は全て排除され、
有限な場合は対称分解と二十印から $Z=X$ である。

二次正規化の場合は定理\ref{thm:quadric-conductor-quintic}の正規化環の表示を用いる。
$A(x)$ の階数三以下の台は四行の反対根の階数一軌跡、
階数二以下の $A(x)$ は正規化繊維長四以上となるため存在しない。
実際四行根の十成分は全二次式の基底で階数零点はない。
定理\ref{thm:rank-two-surface-excess-curves}と\ref{thm:rank-one-total-five}により
その台に曲面はなく曲線の被約総次数は五以下である。
いまの定理で $A$ は有限な五次対称表示となり、定理\ref{thm:finite-quintic-web-exclusion}が非恒等根を排除する。
\end{proof}
\path{computations/verify_quintic_web_curve_excess.py} は複体のEuler多項式と全ての台の数値境界を検算する。
